\documentclass[12pt,a4paper]{article}
\usepackage[utf8]{inputenc}
\usepackage{graphicx}
\usepackage{amsmath}
\usepackage{hyperref}
\usepackage{geometry}
%\geometry{margin=1in}
\usepackage{ebgaramond}
% Packages for bibliography
\usepackage[numbers]{natbib}
% Packages for tables
\usepackage{booktabs}    % For professional tables with better spacing
\usepackage{multirow}    % For merged rows in tables
\usepackage{siunitx}     % For number formatting
\usepackage{array}       % For better column formatting
\usepackage{float}
% Packages for bibliography
\usepackage[numbers]{natbib}

% Define title and author
\title{EnvisionHGdetector: Skeleton-based Co-speech Gesture Detection Using Mediapipe and a ConvNet Classifier}
\author{Wim Pouw (1, 3) \and Bosco Yung \and Sharjeel Shaikh \and James Trujillo (3,4) \and Gerard de Melo  \and Babajide Owoyele (2, 3)}
\date{}

\begin{document}
\maketitle

% Author Note
\section*{Author Note}
1. Donders Institute for Brain, Cognition and Behaviour, Radboud University, Nijmegen, Netherlands.\\
2. Hasso Plattner Institute, University of Potsdam, Germany.\\
3. for the EnvisionBOX team.\\
4. University of Amsterdam, Netherlands. \\

Acknowledgments: WP is funded by a DFG VICOM grant (PO 2841/1-1), a NWO VENI grant (VI.Veni 0.201G.047), and WP, GM, BO, are supported by a NWO XS grant (406.XS.24.01.043). \\
Contact Information: Wim Pouw (email: wim.pouw@donders.ru.nl), Babajide Owoyele (email: Babajide.Owoyele@hpi.de)

\newpage

\begin{abstract}
We introduce \texttt{envisionHGdetector}, an open-source Python package for automatic detection of co-speech gestures and self-adaptors in single-person videos. Unlike existing approaches that require custom training, our tool works out-of-the-box on frontal view recordings through a convolutional neural network trained on over 7,000 annotated instances from diverse datasets (Zhubo, SAGA, Multisimo, and TEDM3D). Building on NoddingPigeon \cite{yungNoddingPigeon2022}, our CPU-efficient implementation combines MediaPipe Holistic skeleton tracking with a lightweight CNN classifier, achieving 78.4\% accuracy in gesture detection (FPR = .213, FNR = .512). The package outputs gesture segment labels, confidence scores, segmented videos, and ELAN files. To facilitate further research, we provide the source code (https://github.com/WimPouw/envisionhgdetector), training features (https://osf.io/qgecj/), and plan to release the \textit{envisionGestureDataset}, a deidentified machine-readable dataset for training. This tool aims to augment human annotation in gesture studies, providing a general-purpose solution for researchers studying multimodal communication.

\noindent\textbf{Keywords:} Co-speech gesture detection, MediaPipe, machine learning, multimodal communication, gesture studies
\end{abstract}

\newpage

\section{Introduction}
Despite an increasing interest in multimodal aspects of language production, to date there is no directly applicable co-speech gesture detector designed to work on a variety of recordings of speaking subjects. The lack of such a tool is limiting our ability to study gesture on a grand scale, which could accelerate research of human communication both in the wild and in the lab.

Many researchers are convinced that a better understanding of human hand gestures will provide deeper insights into how humans think and what the ultimate format is of their communication \citep{feyereisenCognitivePsychologySpeechrelated2017, kendonGestureVisibleAction2004, mcneillHandMindWhat1992, goldin-meadowGestureSignLanguage2017}. For linguists specifically, understanding gesture means understanding how well-formed aspects of languages are entangled with more imagistic, situated, and pragmatic modes of conveying meaning using a variety of gestural forms \cite{alviarMultimodalCoordinationPragmatic2023, murgianoSituatingLanguageRealworld2021}. Such forms include beat gestures moving rhythmically with speech prosody, iconic and metaphoric gestures that shape to represent something visuospatially \citep{streeckDepictingGesture2008}, and other forms like pointing, pragmatic gestures for turn-taking, and emblematic gestures.

To know what the difference is between scratching your head and gesturing to scratch your head to signal ignorance is to know a distinction that so far only humans really have been able to make. While this often requires broader context, in many cases the forms that hand movements make when gesturing are quite differentiable from other types of movements. Gestures have their own unique rhythms and postures \cite{pouwEntrainmentModulationGesture2019, streeckDepictingGesture2008, morgensternTheeSingOpening2021}, share commonalities within cultures \cite{kendonGestureVisibleAction2004}, and may diverge from other types of actions to maintain clear meaning. In the current era of deep learning \citep{sieversDeepSocialNeuroscience2024}, a machine classifier should in principle do relatively well in differentiating gestures from non-gesture movements.

In this paper, we introduce the Python package envisionHGdetector, a light-weight tool for classifying co-speech gestures from video as well as other actions such as self-adaptors. Additionally, we raise a challenge to the community: Can we pool our data collection resources to develop better machine learning for a truly versatile social action recognition system? To address this challenge, we plan to release a machine-readable deidentified gesture dataset combining multiple sources that could serve as a benchmark training dataset.

\section{Related Work}
Current tools for multimodal analysis include OpenFace for facial gesture detection \citep{baltrusaitisOpenFaceOpenSource2016}, OpenPose \citep{caoOpenPoseRealtimeMultiPerson2021} and MediaPipe \citep{caoOpenPoseRealtimeMultiPerson2021} for pose tracking, and OpenSmile \citep{eybenOpensmileMunichVersatile2010} for speech acoustic feature extraction. However, machine learning methods for gesture detection typically need custom training for each unique dataset \citep{gregoriRoadmapTechnologicalInnovation2023, henleinOutlookAIInnovation2024}.

Research in machine learning has shown promising results in separating gestures from non-gestures. Rule-based approaches like SPUDNIG \citep{ripperdaSpeedingDetectionNoniconic2020} use pose-tracking and speed timeseries analysis, though these methods often have high false-positive rates of 70\% \citep{ripperdaSpeedingDetectionNoniconic2020} to 86\% \citep{pouwSemanticallyRelatedGestures2021}. More sophisticated approaches include Ghaleb et al.'s \citep{ghalebCoSpeechGestureDetection2024} Spatio-Temporal Graph Convolutional Networks for gesture phase detection, achieving 61\% accuracy on gesture strokes. Ienega and colleagues \citep{ienagaSemiautomationGestureAnnotation2022} developed a semi-automated approach using Light Gradient Boosting Machine, reaching 86\% accuracy when trained on just 3\% of hand-annotated data.

The field of human-computer interaction has also made progress in classifying interactive silent stereotyped gestures \citep{molchanovOnlineDetectionClassification2016, frontedduDynamicHandGesture2022}, with models like motionBERT \citep{zhuMotionBERTUnifiedPerspective2023} achieving over 80\% accuracy in differentiating gesture classes. However, these successes have not yet translated to widespread adoption, partly due to limited computer science training among cognitive scientists \cite{sieversDeepSocialNeuroscience2024} and the lack of comprehensive open datasets for benchmark training.

Current annotation practices, which often require frame-by-frame analysis of gesture stroke phases, will not be fully automated in the near future. However, the simpler task of distinguishing gestures from non-gestures shows promise for automation, particularly when combined with human-in-the-loop verification \citep{ripperdaSpeedingDetectionNoniconic2020}. A key challenge remains the generalizability of models across different recording conditions and labeling practices, highlighting the need for comprehensive, openly available training datasets.



\end{document}
